% ---
% id: statistics/thm:binomial-tails
% kind: thm
% title: "Binomial tails in closed form"
% summary: >-
%   The upper and lower tails of a binomial count are regularised incomplete
%   beta integrals, with every boundary case of the count stated. The upper
%   tail is non-decreasing in the success probability.
% prominence: standard
% status: implemented
% uses: []
% assumes: []
% refs: []
% provenance: newly-derived
% implements: certiqs_key.statistics.tails:binomial_cdf
% lean:
% ---

The count of one outcome over independent rounds with one law is binomial,
and the box construction compares its two tails with their budgets. The
theorem writes both tails as regularised incomplete beta integrals, deriving
the normalisation of the beta integral inside the proof.

\begin{theorem}\label{thm:binomial-tails}
Let $n\ge1$, $p\in[0,1]$, and let $\Ncount{}$ be a binomial random variable
with parameters $(n,p)$, i.e.\
$\Pr[\Ncount{}=i]=\binom{n}{i}p^{i}(1-p)^{n-i}$ for $0\le i\le n$. Write
\begin{equation}\label{eq:ireg}
  \Ireg{x}{a}{b}\coloneqq
  \frac{\int_0^{x}t^{a-1}(1-t)^{b-1}\,\mathrm{d}t}
       {\int_0^{1}t^{a-1}(1-t)^{b-1}\,\mathrm{d}t},
  \qquad x\in[0,1],\ a,b\ge1 .
\end{equation}
Then
\begin{align}
  \Pr[\Ncount{}\ge k]&=\frac{n!}{(k-1)!\,(n-k)!}
    \int_0^{p}t^{k-1}(1-t)^{n-k}\,\mathrm{d}t
    =\Ireg{p}{k}{n-k+1}, &&1\le k\le n,
    \label{eq:uppertail}\\
  \Pr[\Ncount{}\le k]&=1-\Pr[\Ncount{}\ge k+1]
    =\Ireg{1-p}{n-k}{k+1}, &&0\le k\le n-1,
    \label{eq:lowertail}
\end{align}
together with the boundary values $\Pr[\Ncount{}\ge k]=1$ for $k\le0$,
$\Pr[\Ncount{}\ge k]=0$ for $k>n$, $\Pr[\Ncount{}\le k]=0$ for $k<0$ and
$\Pr[\Ncount{}\le k]=1$ for $k\ge n$. Moreover $p\mapsto\Pr[\Ncount{}\ge k]$
is non-decreasing on $[0,1]$ for every $k\ge1$.
\end{theorem}

\begin{proof}
\emph{Idea.} The upper tail is a polynomial in $p$ vanishing at $p=0$ whose
derivative telescopes to a single term; integrating that term gives the first
form of \eqref{eq:uppertail}, and its value at $p=1$ gives the normalisation.
The lower tail is the upper tail of $n-\Ncount{}$.

\emph{Step 1 (the derivative of the upper tail).} Fix $k$ with $1\le k\le n$
and define the polynomial
$G(p)\coloneqq\sum_{i=k}^{n}\binom{n}{i}p^{i}(1-p)^{n-i}$, so that
$G(p)=\Pr[\Ncount{}\ge k]$ and $G(0)=0$ (every summand carries $p^{i}$ with
$i\ge k\ge1$). Set
\begin{equation*}
  a_i\coloneqq n\binom{n-1}{i-1}p^{i-1}(1-p)^{n-i},
  \qquad 1\le i\le n ,
\end{equation*}
a polynomial in $p$ for every such $i$: the exponent $n-i$ is non-negative
throughout, so no negative power of $1-p$ occurs anywhere in this proof and no
division is performed. Differentiate the summands of $G$ separately in the two
cases.

For $k\le i\le n-1$, the product rule and the absorption identities
$i\binom{n}{i}=n\binom{n-1}{i-1}$ and $(n-i)\binom{n}{i}=n\binom{n-1}{i}$ give
\begin{equation*}
  \frac{\mathrm{d}}{\mathrm{d}p}\Big[\binom{n}{i}p^{i}(1-p)^{n-i}\Big]
  =i\binom{n}{i}p^{i-1}(1-p)^{n-i}-(n-i)\binom{n}{i}p^{i}(1-p)^{n-i-1}
  =a_i-a_{i+1},
\end{equation*}
where $i\ge k\ge1$ makes the first power non-negative, $i\le n-1$ makes the
second one non-negative, and $a_{i+1}=n\binom{n-1}{i}p^{i}(1-p)^{n-i-1}$ by
substituting $i+1$ for $i$ in the definition of $a_i$, which is legitimate
because $i+1\le n$. For $i=n$ the summand is $p^{n}$, whose derivative is
$np^{n-1}=a_n$; this is the case the general formula would have written with a
factor $(1-p)^{-1}$, and it is treated separately for exactly that reason.
Summing the two cases, the sum telescopes,
\begin{equation*}
  G'(p)=\sum_{i=k}^{n-1}\big(a_i-a_{i+1}\big)+a_n
       =\big(a_k-a_n\big)+a_n
       =a_k ,
\end{equation*}
which is also correct at $k=n$, where the telescoping sum is empty and
$G'(p)=a_n=a_k$.

\emph{Step 2 (the integral form and its normalisation).} Spelling $a_k$ out,
$G'(p)=n\binom{n-1}{k-1}p^{k-1}(1-p)^{n-k}$, and by the fundamental
theorem of calculus with $G(0)=0$,
\begin{equation}\label{eq:Gint}
  G(p)=n\binom{n-1}{k-1}\int_0^{p}t^{k-1}(1-t)^{n-k}\,\mathrm{d}t,
  \qquad
  n\binom{n-1}{k-1}=\frac{n!}{(k-1)!\,(n-k)!}.
\end{equation}
This is the first equality of \eqref{eq:uppertail}. For the second, evaluate
\eqref{eq:Gint} at $p=1$: there $\Ncount{}=n$ with probability one and $k\le
n$, so $G(1)=1$, whence
$\int_0^{1}t^{k-1}(1-t)^{n-k}\,\mathrm{d}t=\frac{(k-1)!\,(n-k)!}{n!}$;
dividing \eqref{eq:Gint} by this value gives exactly
$\Ireg{p}{k}{n-k+1}$ as defined in \eqref{eq:ireg}, with $a=k\ge1$ and
$b=n-k+1\ge1$.

\emph{Step 3 (the lower tail).} For \eqref{eq:lowertail}, note first that
$\{\Ncount{}\le k\}$ and $\{\Ncount{}\ge k+1\}$ partition the sample space.
Next, $n-\Ncount{}$ is binomial with parameters $(n,1-p)$, because
$\Pr[n-\Ncount{}=i]=\Pr[\Ncount{}=n-i]
=\binom{n}{n-i}p^{n-i}(1-p)^{i}=\binom{n}{i}(1-p)^{i}p^{n-i}$. Hence, for
$0\le k\le n-1$ (so that $n-k\ge1$),
$\Pr[\Ncount{}\le k]=\Pr[n-\Ncount{}\ge n-k]=\Ireg{1-p}{n-k}{k+1}$ by
\eqref{eq:uppertail} applied to $n-\Ncount{}$.

\emph{Step 4 (boundary values and monotonicity).} The boundary values are
immediate from $0\le\Ncount{}\le n$. Monotonicity in $p$ follows from
$G'(p)=n\binom{n-1}{k-1}p^{k-1}(1-p)^{n-k}\ge0$ on $[0,1]$ for $1\le k\le n$;
for $k>n$ the tail is the constant $0$ by the boundary values.
\end{proof}
